\section{Resultados y plan de análisis}

\subsection{Verificación disponible}

La auditoría reproducible encuentra 8.000/1.490/17.194 imágenes en
entrenamiento/validación/test, coincidiendo con BMAD. Una muestra estratificada
de 308 ficheros se abrió correctamente y todos fueron PNG monocromos de
1024$\times$1024. Las pruebas del software cubren carga de ambas clases, forma de
salida de los tres modelos, empates en AUROC y selección del umbral.

Se completó primero una prueba técnica de una época con 32 imágenes normales y
16 por clase. Después se ejecutó un \textbf{piloto} más exigente: tres épocas,
las 8.000 imágenes de entrenamiento, validación y test completos y semilla 42.
La tabla resume ese piloto; no sustituye las tres semillas y 20 épocas del
protocolo definitivo.

\begin{center}
\begin{tabular}{lrrrrr}
\toprule
Modelo & Parámetros & AUROC & AUPRC & Bal. acc. & Tiempo (s) \\
\midrule
AE       & 609.857 & 0,4818 & 0,9552 & 0,5180 & 581,7 \\
VAE      & 740.993 & 0,4934 & 0,9567 & 0,5218 & 592,4 \\
GANomaly & 959.586 & \textbf{0,6019} & \textbf{0,9691} & \textbf{0,5865} & 610,4 \\
\bottomrule
\end{tabular}
\end{center}

La prevalencia anómala de test es 95,46\%; por eso AUPRC ronda 0,95 incluso en
AE y VAE, cuyos AUROC están cerca del azar. El umbral de validación produce
sensibilidad/especificidad 0,179/0,857 en AE, 0,215/0,828 en VAE y 0,475/0,698
en GANomaly. La alternativa adversarial ofrece la única señal prometedora, pero
una sola semilla y tres épocas no permiten atribuir una ventaja estable.

\subsection{Resultado científico pendiente de cómputo}

La tabla definitiva se generará automáticamente tras ejecutar 20 épocas sobre
las 8.000 imágenes y tres semillas. Debe contener media $\pm$ desviación de
AUROC, AUPRC y exactitud balanceada, además del tiempo. No se rellenan cifras
antes de ejecutar el protocolo completo. La hipótesis $H_1$ es que al menos VAE
o GANomaly mejora el AUROC medio del AE; $H_0$ afirma que la diferencia no es
consistente entre semillas.

El análisis no se limitará al mejor número: se inspeccionarán distribuciones de
puntuación, falsos positivos/negativos y reconstrucciones. Una ganancia pequeña
con gran coste o variabilidad no justificará declarar superior el método más
complejo. También se comparará el ranking con controles simples de intensidad
para detectar atajos no anatómicos.
